\section{Overview of UniTime and CPSolver}
\label{sec:unitime}

\subsection{The UniTime system}

UniTime is a comprehensive, open-source academic scheduling system released under the Apache
License, Version 2.0, and maintained as a sponsored project of the Apereo Foundation. It is
web-based, implemented in Java, and deployed at large institutions worldwide, including Purdue
University~\cite{rudova2011complex}. Later accounts of the same system cover elective courses and
multiple sections~\cite{muller2016curriculum}. The system comprises four integrated components: course timetabling, examination
timetabling, student scheduling, and event management. Only the first is used in this project,
although the data model shared by all four determines the structure of the import files described
in Section~\ref{sec:pipeline}.

Within UniTime a course is modelled hierarchically. A \emph{course offering} belongs to an
\emph{instructional offering}, which carries one or more \emph{configurations}; a configuration
contains \emph{scheduling subparts} (for example a lecture subpart and a laboratory subpart); and
each subpart contains the individual \emph{classes} that are actually placed in the timetable. A
class inherits its permitted meeting patterns from its subpart. This hierarchy is what allows a
course with joint lectures and parallel laboratory batches to be represented as a single offering,
and it is the reason the institution's division-and-batch structure maps onto UniTime without
distortion~\cite{murray2007modeling}.

\subsection{The constraint solver}

The embedded solver is CPSolver~\cite{cpsolver}, a local-search framework with constraint-programming primitives,
distributed under the GNU LGPL. It won two tracks of the International Timetabling Competition 2007
and reached the final in all three~\cite{muller2009itc,muller2025itc}. In the course-timetabling model the variable is a class, the
value is a placement of time and room, and the objective is a weighted sum of soft-constraint
criteria.

\subsection{Iterative Forward Search}

The solver is built on Iterative Forward Search (IFS)~\cite{muller2004}. IFS belongs to the family
of local-search methods, but differs from classical local search in that it operates over feasible
yet potentially \emph{incomplete} solutions: some classes may remain unassigned, but every hard
constraint holds for the classes that are assigned. Intermediate solutions are therefore always
interpretable and actionable by a human scheduler, in contrast to complete-but-infeasible solutions
produced by unconstrained search.

From an initial timetable $\omega$, each iteration selects a class and a placement, removes the
assignments that hard-conflict with it, and installs the new one. The best timetable seen,
$\sigma$, is kept and returned at the end, so the search can be stopped and resumed from any
feasible timetable. Listing~\ref{lst:ifs} states the algorithm; Figure~\ref{fig:solverloop} is the
same loop as it actually runs, including the distinction between proposing a move and committing
it. The learning layer of Section~\ref{sec:rl} replaces only the \emph{Propose move} box. That step
draws one assignment from the neighbourhood of the current one~\cite{muller2005thesis}. Later
changes to the competition solver's time changes, room changes and swaps are described
in~\cite{muller2016evolution}.

\begin{lstlisting}[style=plainlst,caption={Iterative Forward Search},label={lst:ifs}]
function IFS(omega):
    sigma = omega
    while canContinue(omega) do:
        v     = selectVariable(omega)            // select a class
        d     = selectValue(omega, v)            // select a placement
        gamma = hardConflicts(omega, v/d)        // conflicting assignments
        omega = (omega \ gamma) union {v/d}      // remove conflicts, assign
        if better(omega, sigma) then
            sigma = omega                        // update best solution
    end while
    return sigma
end function
\end{lstlisting}

\begin{figure}[H]
\centering
\begin{tikzpicture}[
  node distance=5.2mm and 8mm,
  box/.style={draw,rounded corners=2pt,align=center,font=\small,minimum height=8mm,inner xsep=5pt,fill=coepblue!5},
  dec/.style={draw,diamond,aspect=2.1,align=center,font=\scriptsize,inner sep=1pt,fill=coepblue!10},
  term/.style={draw,rounded corners=8pt,align=center,font=\small,fill=gray!12},
  ar/.style={-{Stealth[length=2mm]}}
]
\node[term] (start) {Start / resume};
\node[dec, below=of start] (stop) {Stop\\check};
\node[box, below=of stop] (prop) {Propose move\\{\scriptsize(select class, select placement)}};
\node[dec, below=of prop] (sel) {Move\\selected?};
\node[box, below=of sel] (apply) {Apply move\\{\scriptsize remove hard conflicts, assign}};
\node[box, below=of apply] (score) {Score solution\\{\scriptsize weighted soft criteria}};
\node[dec, below=of score] (better) {Is better\\than best?};
\node[box, below=of better] (save) {Save best $\sigma$};
\node[term, right=22mm of stop] (end) {Return $\sigma$};
\node[box, right=14mm of sel, text width=26mm] (adv) {Advance counter\\{\scriptsize (no commit)}};

\draw[ar] (start) -- (stop);
\draw[ar] (stop) -- node[left,font=\scriptsize]{continue} (prop);
\draw[ar] (stop) -- node[above,font=\scriptsize]{timeout} (end);
\draw[ar] (prop) -- (sel);
\draw[ar] (sel) -- node[left,font=\scriptsize]{yes} (apply);
\draw[ar] (sel) -- node[above,font=\scriptsize]{no} (adv);
\draw[ar] (apply) -- (score);
\draw[ar] (score) -- (better);
\draw[ar] (better) -- node[left,font=\scriptsize]{yes} (save);

% feedback edges
\draw[ar] (adv.north) |- ($(prop.north)+(0,5mm)$) -| (prop.north);
\draw[ar] (better.east) -- node[above,font=\scriptsize]{no} ++(34mm,0) |- ($(stop.south)+(0,-2.5mm)$);
\draw[ar] (save.west) -- ++(-20mm,0) |- (stop.west);
\end{tikzpicture}
\caption{The CPSolver iteration loop. The incumbent $\sigma$ is updated only when the comparator
judges the new solution better. Section~\ref{sec:rl} substitutes only the \emph{Propose move} step.}
\label{fig:solverloop}
\end{figure}

\subsection{Conflict-based Statistics}

To avoid cycling and to improve final quality, IFS is augmented with Conflict-based Statistics
(CBS)~\cite{muller2004}. CBS records how often each assignment has historically caused a hard
conflict with another:
\[
  \mathrm{CBS}\bigl[A/a \rightarrow \lnot B/b\bigr] = c_{ab},
\]
meaning that assigning class $A$ the placement $a$ has unassigned $B/b$ on $c_{ab}$ previous
occasions. During value selection each candidate placement is penalised by the accumulated conflict
frequency of the assignments it would displace, steering the search away from historically
problematic regions. When no complete solution is found, the contents of the CBS structure identify
the hard constraints responsible for the infeasibility, which can be presented to the user as
guidance for relaxation. In the configuration used throughout this project CBS is enabled
(\texttt{General.CBS=true}) with the conflict-statistics extension registered on the solver.

\subsection{Constraint enforcement guarantees}

The framework enforces the invariant that hard constraints are never violated by an assignment that
persists in the current timetable; infeasible placements are removed rather than tolerated. Arc
consistency can in principle be maintained throughout the search, and IFS can be made complete
through dynamic backtracking, but maintaining full arc consistency is impractical at realistic
instance sizes, so the heuristic form is used. In the configuration used here
\texttt{Basic.DisobeyHard} is \texttt{false}, which means that every mechanism described in
Section~\ref{sec:rl} inherits this guarantee without needing its own hard-constraint checking.
